\section{Nghiên cứu liên quan}
\label{sec:related_work}

Ghép ảnh 2D gia tăng từ video máy bay không người lái (UAV) đòi hỏi giải quyết đồng thời bốn thách thức kỹ thuật: đăng ký chuỗi khung hình, giới hạn tài nguyên tính toán, phát hiện rủi ro lỗi đăng ký và kiểm soát lan truyền sai số. Phần này hệ thống hoá các công trình tiền nhiệm theo bốn hướng chính và xác lập phạm vi so sánh thực nghiệm của nghiên cứu.

\subsection{Ghép ảnh UAV gia tăng và SLAM đơn mục}
Các hệ thống ghép video UAV truyền thống và đương đại chủ yếu dựa trên luồng ước lượng xạ ảnh (homography) tuần tự. Hwang cùng cộng sự~\cite{hwang2026} phát triển hệ thống ghép ảnh trực giao 2D thời gian thực trên CPU, ước lượng ma trận homography gia tăng bằng các bộ trích xuất đặc trưng cổ điển (SIFT~\cite{lowe2004}, ORB~\cite{rublee2011}) kết hợp RANSAC~\cite{fischler1981}, sử dụng ngưỡng số lượng inlier cố định để kiểm duyệt biến đổi giữa các khung hình liên tiếp. Li cùng cộng sự~\cite{li2023} đề xuất khung ghép video gia tăng trong viễn thám UAV, ứng dụng tối ưu hoá sai số chiếu lại có trọng số, chọn khung hình chính (keyframe) theo tỷ lệ đè phủ quang học và cơ chế kết nối lại (rematch) với keyframe lịch sử nhằm giảm thiểu trôi dạt theo dõi. Chen cùng cộng sự~\cite{chen2022} xử lý biến dạng phối cảnh khi độ cao bay thay đổi thông qua quy hoạch động tìm đường nối tối ưu kết hợp phép biến đổi nửa xạ ảnh bảo toàn hình dạng. Nhằm khắc phục tính dễ tổn thương của chuỗi homography tuần tự, Bu cùng cộng sự~\cite{map2dfusion} giới thiệu Map2DFusion, kết hợp theo dõi trực tiếp mức pixel từ SLAM thị giác đơn mục với đồ thị keyframe và tối ưu hoá bundle adjustment cục bộ. Khác với các hệ thống SLAM thị giác toàn diện duy trì điểm bản đồ 3D và tối ưu hoá bundle adjustment, nghiên cứu của chúng tôi tập trung vào pipeline ghép ảnh 2D gia tăng bằng homography nhẹ vận hành dưới ràng buộc tài nguyên tuần tự mà không sử dụng tối ưu hoá đồ thị toàn cục.

\subsection{Đối sánh đặc trưng học sâu và ràng buộc tài nguyên}
Học sâu đã mang lại những bước tiến lớn trong đối sánh đặc trưng dưới sự thay đổi mạnh về điều kiện chiếu sáng và góc nhìn. DeTone cùng cộng sự~\cite{detone2018} giới thiệu SuperPoint, mạng tự giám sát để đồng thời phát hiện điểm đặc trưng và trích xuất vector mô tả. Lindenberger cùng cộng sự~\cite{lindenberger2023} phát triển LightGlue, mạng nơ-ron đồ thị tích hợp cơ chế loại bỏ điểm không khớp thích ứng và dừng sớm để tiết kiệm chi phí. Đối với hướng tiếp cận không dùng bộ dò điểm (detector-free), Sun cùng cộng sự~\cite{sun2021} xây dựng LoFTR dựa trên kiến trúc Transformer coarse-to-fine, trong khi Wang cùng cộng sự~\cite{wang2024} đề xuất Efficient LoFTR nhằm giảm độ phức tạp tính toán thông qua cơ chế chú ý tổng hợp (aggregated attention) với chọn lọc token thích ứng và tầng tương quan hai giai đoạn. Nhằm phục vụ các nền tảng biên, Potje cùng cộng sự~\cite{xfeat2024} thiết kế XFeat, kiến trúc đối sánh nhẹ tối ưu hoá riêng cho thiết bị bị giới hạn tài nguyên. Gần đây, Kim và Kim~\cite{kim2025} kết hợp trích xuất SIFT với bộ đối sánh LightGlue để xử lý các cặp ảnh UAV có góc xoay lớn.

Mặc dù các mô hình học sâu thể hiện năng lực đăng ký cặp ấn tượng, hành vi của chúng trong chuỗi video UAV liên tục bộc lộ những đánh đổi thực tế. Trong các chuỗi bay dài thiếu vắng bundle adjustment toàn cục, các biến đổi học sâu tuần tự vẫn có thể tích lũy sai số subpixel qua các bước bay danh định. Điều này đặt ra yêu cầu khảo sát các chiến lược kích hoạt có chọn lọc thay vì gọi mô hình sâu liên tục trên mọi bước.

\subsection{Dự báo lỗi đăng ký và ước lượng độ bất định}
Zaffar cùng cộng sự~\cite{zaffar2024} so sánh độ bất định dựa trên truy hồi, độ bất định aleatoric học từ dữ liệu và kiểm chứng hình học trong nhận diện địa điểm; đồng thời đề xuất một đường cơ sở sử dụng tư thế của ảnh tham chiếu. Mục tiêu của họ là độ tin cậy của kết quả truy hồi ảnh, không phải drift tích lũy của mosaic UAV; vì vậy cần đánh giá chỉ báo bằng một thước đo thay thế lỗi đăng ký được định nghĩa rõ. Barath cùng cộng sự~\cite{barath2023} cung cấp homography tham chiếu để đánh giá bộ ước lượng bền vững và phương pháp lọc tương ứng điểm. Ivashechkin cùng cộng sự~\cite{ivashechkin2021} giới thiệu VSAC, sử dụng khái niệm inlier độc lập và kiểm tra giả thuyết tuần tự thích ứng để cải thiện ước lượng bền vững. Các công trình này phân biệt đánh giá độ tin cậy và ước lượng bền vững với bài toán quản lý một phép biến đổi đã được chấp nhận trong chuỗi dài.

Đầu ra nội tại của bộ ước lượng là tín hiệu hữu ích khi chấp nhận phép đăng ký~\cite{hwang2026,li2023}, nhưng phần dư nhỏ trên tập correspondence được chọn không tự chứng minh tính nhất quán toàn cục. Nghiên cứu này khảo sát liệu tính nhất quán diện mạo hậu đăng ký và chỉ báo hình học có bổ sung bằng chứng mô tả ngoài số điểm hay không; không giả định nhất quán trắc quang chứng minh hình học đúng.

\subsection{Lan truyền sai số chuỗi và chính sách phục hồi}
Khi xảy ra suy giảm chất lượng đăng ký, hệ thống ghép ảnh gia tăng cần một chính sách phục hồi để xử lý bước lỗi. Giải pháp đơn giản nhất là loại bỏ khung hình (frame dropping) kém tin cậy~\cite{hwang2026}. Tuy nhiên, việc loại bỏ khung hình làm nới rộng khoảng cách thời gian giữa khung hình mới và điểm neo hiện thời của bản đồ, làm tăng độ lệch phối cảnh và có thể dẫn đến các bước loại bỏ tiếp theo. Li cùng cộng sự~\cite{li2023} giải quyết vấn đề này bằng cách tìm kiếm các keyframe trong quá khứ để đối sánh lại, dù khả năng kết nối lại phụ thuộc vào độ đè phủ không gian và giới hạn cửa sổ tìm kiếm. Kharismawati và Kazic~\cite{kharismawati2025} phát triển pipeline DroneZaic cho video UAV nông nghiệp, sử dụng phân cụm chuyển động để chia chuyến bay thành các mảnh độc lập (submosaics).

Chiến lược tách mảnh bảo toàn toàn bộ số lượng khung hình gốc bằng cách mở điểm neo tham chiếu mới khi gặp đứt gãy đăng ký, song bản chất là đánh đổi tính liên tục của bề mặt bản đồ đơn nhất lấy độ bao phủ cục bộ. Ngược lại, chiến lược loại bỏ khung hình duy trì tính liền mạch của chuỗi sống sót nhưng có nguy cơ loại bỏ phần lớn hành trình bay khi tình trạng suy giảm kéo dài. Trong nghiên cứu này, chúng tôi so sánh các chính sách phục hồi trên dưới các tiêu chí kích hoạt rủi ro thống nhất.
